\documentclass[11pt]{article}
\let\origbibliographystyle\bibliographystyle
\renewcommand{\bibliographystyle}[1]{}
\usepackage[preprint]{acl}
\let\bibliographystyle\origbibliographystyle
\usepackage{graphicx}
\usepackage{booktabs}
\usepackage{amsmath,amssymb}
\usepackage{xcolor}
\setcitestyle{numbers,square}

\newcommand{\fillin}[1]{\textcolor{red}{[FILL IN: #1]}}
\newcommand{\rec}[1]{\textcolor{blue}{[SUGGESTED: #1]}}
% Colour key: red = still to fill in; blue = suggestion not yet acted on;
% green = text updated on 19 Nov to match the implemented pipeline;
% magenta = cannot be written until results exist.
\definecolor{updgreen}{rgb}{0,0.5,0}
\newcommand{\upd}[1]{{\color{updgreen}#1}}
\newcommand{\pending}[1]{\textcolor{magenta}{[PENDING RESULTS: #1]}}

\title{Switches or Dials? How Language Model Traits Respond to Activation Steering}
\author{Lauren Shen \and Mahdi Afshari \and Katelyn Gan \and Danial Hosseintabar \\
Massachusetts Institute of Technology}
\date{}

\begin{document}
\maketitle

\begin{abstract}
\fillin{4--6 sentence abstract.}
\upd{Activation steering changes a language model's behavior by adding a trait
direction to its activations, but how behavior responds to the steering strength
is unknown. Under a dial, each reply gradually becomes more aligned with the
trait as strength grows; under a switch, each prompt flips abruptly at its own
threshold. The question is still open because the smooth curves reported in prior
work may occur only because measurements average over different prompts. We
measure, for ten traits across five Qwen3 model sizes with Gemma 3 as a second
family, how much of the variation in steering response lies within prompts
versus between them. We fit a hierarchical probit model to individual replies and
report the switchiness $S$, the share of the population spread that comes from
prompts differing in threshold. The full pipeline is implemented, and its
estimator recovers known switch and dial settings in simulation.}
\pending{one sentence on the main finding: $S$ per trait group on Qwen3-4B.}
\end{abstract}

%------------------------------------------------------------
\section{Introduction}

Language models might be the most fascinating invention of the decade but they often
hallucinate facts or refuse to answer prompts. To influence the model, you may use a
technique called activation steering to affect a specific trait \cite{chen2025persona}
\upd{by adding a multiple of a trait direction, named the steering vector, to the
model's activations during generation.}

There are currently two hypotheses on how increasing the strength of these steering
vectors impacts the model: dial and switch. The dial represents how, as the magnitude
of the steering vector is incremented, each generated response gradually becomes
aligned with the trait represented by the steering vector. On the other hand, the
switch represents how each prompt flips abruptly from normal to trait-like behavior at
a specific threshold. This question is still open because the smooth curves reported
in prior work may occur because measurements average over different prompts, use a
handful of strengths, and study a single model.

If steering works like a dial, a trait's strength chosen on test prompts will have a
similar, predictable effect on new prompts. If it works like a switch, the same
strength will fully change some new prompts and not affect others at all, so safety
tools built on steering can fail without warning. \upd{Today strength is chosen by
sweeping, that is, trying a range of strengths and keeping the one that works best
on average, although the same strength is known to work well on some inputs and
poorly on others \cite{braun2025unreliability}. Recent studies also disagree on
whether steering along random directions causes side effects
\cite{malla2026interference,ong2026forecasting}.}

\fillin{A paragraph stating what this paper does. The updated proposal dropped it.}
\upd{We measure, for ten traits across five model sizes, how much of the variation
in steering response lies within prompts versus between them, which kinds of
traits act as switches, when during generation the switch is decided, and how all
of this changes with scale.}

\fillin{Explicit list of research questions and contributions.}
\upd{Our research questions are:
\begin{itemize}
\item \textbf{Q1, switches or dials?} How much of the spread of the population
response curve comes from prompts having different thresholds, and how much from
each prompt responding gradually?
\item \textbf{Q2, which traits switch?} Do decision traits switch more sharply than
style traits?
\item \textbf{Q3, where is the switch decided?} Is a reply's mode fixed by the
prompt representation, by its first few tokens, or throughout generation?
\item \textbf{Q4, regimes and scale.} How do the threshold, the switchiness, and the
usable strength window change with model size, and do side effects have two
strength regimes?
\end{itemize}
Our contributions are the switchiness statistic $S$ with an estimator that
separates within-prompt from between-prompt variation, position-restricted
steering interventions for Q3, and an open pipeline that runs entirely on local
models.}

\fillin{One sentence on significance to NLP research and to a real-world problem
(rubric item).}
\upd{For NLP research, $S$ turns a qualitative disagreement about steering into a
measurable quantity; in practice, if traits turn out to be switches, steering
strength should be calibrated per prompt and not globally, and $S$ gives a single
number for deciding when that is necessary.}

%------------------------------------------------------------
\section{Related Work}

\fillin{1--2 sentences of background on activation steering for readers outside
interpretability (residual stream, steering vector, difference-of-means).}
\upd{A transformer's residual stream is the vector, one per token position, that
each layer reads from and adds to. A steering vector is a direction in that space
associated with a behavior; a common way to obtain one is the difference of
means, which subtracts the average activation over replies that lack the behavior
from the average over replies that show it \cite{chen2025persona}. Adding a
multiple of the vector during generation shifts the model toward the behavior.}

\subsection{Trait Geometry}
Much of the related work studies how these traits are arranged relative to one another
inside the model, where each trait is represented as a direction
\cite{lu2026assistant,xu2026manifold,abootorabi2026geometry}.
\fillin{One sentence on how your project differs from this line of work.}
\upd{These works describe where traits sit; we ask how behavior moves when a model
is pushed along those directions.}

\subsection{Steering Strength and Reliability}
Braun et al.\ \cite{braun2025unreliability} found that the same steering can work well
on some inputs and poorly on others, and Fan et al.\ \cite{fan2026steerable} predict,
from the model's early internal state, whether a given prompt will be pushed too
little, the right amount, or too much. Both treat differences between prompts as noise
to predict. Neither asks whether that variation comes from replies to the same prompt
or from differences between prompts, which is exactly what separates a switch from a
dial.

\subsection{Side Effects}
While Malla et al.\ \cite{malla2026interference} and Ong et al.\
\cite{ong2026forecasting} measure large cross-behavior side-effect matrices, they reach
conflicting conclusions about random directions. Rather than having to re-measure full
matrices, we measure a reduced one to test whether the conflict is explained by two
strength regimes (Q4).

%------------------------------------------------------------
\section{Methodology}

\subsection{Steering Setup}
We \upd{extract} difference-of-means steering vectors for \upd{ten} traits with
the persona-vectors recipe \cite{chen2025persona}, from \upd{three} independently generated
prompt sets per trait. Vectors are extracted and applied at the middle layer of each
model. Strength $a$ is measured in units of the median residual-stream norm at the
steered layer, so that it is comparable across model sizes. An LLM judge will score
every reply \upd{for the steered trait} and for coherence on a 0 to 9 scale, and a reply will
count as expressing the trait when its score is at least a 5. \upd{Scoring every reply
against all ten rubrics would take about one million judge calls per model and
vector set, so all ten rubrics are applied only to the reduced side-effect runs
of Q4.}

\fillin{The final list of traits and how many prompt sets per trait.}
\upd{The decision traits are refusing, sycophantic (agreeing with a false claim),
hallucinating (fabricating), deceptive, and clarifying (asking a clarifying
question); the style traits are humorous, optimistic, poetic, neurotic, and
impolite. Each trait has three extraction sets (A, B, C), each with five pairs of
positive and negative system instructions and 20 questions, giving three vectors
per trait and model. For every pair of instruction and question we sample two
replies under the positive and two under the negative instruction, keep pairs in
which the positive reply expresses the trait, the negative one does not, and both
are coherent, and subtract the mean residual over reply tokens. The vector is
normalised to unit length, and steering adds $a$ times the median residual norm,
measured on 16 neutral prompts. By default it is added at every position from
which a reply token is predicted, that is, from the last prompt token onward.}
\fillin{Why the middle layer (rubric: justify design choices).}
\upd{We use block $L/2$ of $L$ in every model. A rule fixed in advance gives the
same relative depth at every model size, and it avoids choosing a layer per trait
by its steering effect, which would tune the very response we measure. Middle
layers are where the persona-vectors work reports steering to be most effective
\cite{chen2025persona}.}
\fillin{A figure giving an overview of the pipeline or a schematic of dial vs.\
switch response curves.}

\subsection{Q1: Switches or Dials?}
For prompt $i$, let $p_i(a)$ be the probability that a sampled reply expresses the
trait. We model
\begin{equation}
p_i(a) = \Phi\!\left(\frac{\log a - \log t_i}{w}\right),
\end{equation}
where $\Phi$ is the standard normal cumulative distribution, $t_i$ is the prompt's own
threshold, and $w$ is the width of its transition. Thresholds vary across prompts as
$\log t_i \sim \mathcal{N}(\mu, \sigma_t^2)$. Averaging over prompts then gives the
population curve
\begin{equation}
\Phi\!\left(\frac{\log a - \mu}{\sqrt{w^2 + \sigma_t^2}}\right),
\end{equation}
so a population-level fit cannot separate $w$ from $\sigma_t$.

Therefore, we fit the hierarchical model to individual replies, where repeated samples
of the same prompt identify $w$, and report the switchiness
\begin{equation}
S = \frac{\sigma_t^2}{\sigma_t^2 + w^2},
\end{equation}
which is the share of the population spread that comes from prompts differing in
threshold. $S$ near 1 is a switch, while $S$ near 0 is a dial in which every prompt
responds gradually.

\upd{Equation 1 forces $p_i \to 0$ as $a \to 0$, but some replies express a trait
without steering, so the fitted model adds a per-trait baseline rate $g$,
identified by the unsteered replies: $p_i(a) = g + (1-g)\,\Phi\big((\log a - \log
t_i)/w\big)$. We fit only strengths below the trait's coherence break
(Section~\ref{sec:q4}), because beyond it trait scores fall again as replies turn
to noise, which a monotone model cannot represent. As a reply-level form of the
same question, we also report whether individual judge scores are bimodal at the
strength where about half of the replies express the trait (Sarle's bimodality
coefficient).}

\fillin{How the hierarchical model is fit (estimation method, software).}
\upd{The data for each trait are counts: for prompt $i$ and strength $a$, how many
of the 8 sampled replies express the trait. We estimate $\mu$, $\sigma_t$, $w$,
and $g$ by marginal maximum likelihood. The prompt threshold is integrated out
numerically on a fixed grid of 181 points, and the likelihood is maximised with
L-BFGS-B from nine starting points (our own implementation in NumPy and SciPy).
Widths are bounded below at 0.02 in log strength; a fitted $w$ on that bound
means the transition is narrower than the strength grid can resolve, and we
report it as such. Intervals are 95\% percentile intervals from 200 resamples of
prompts.}

\subsection{Q2: Which Traits Switch?}
We hypothesize that decision traits, where a reply either takes an action or not
(refusing, agreeing with a false claim, fabricating, deceiving, or asking a clarifying
question), have higher $S$ than style traits that can be expressed by degree (humor,
optimism, poetic or neurotic tone). We assign traits to groups before measuring them.

\subsection{Q3: Where Is the Switch Decided?}
Fan et al.\ \cite{fan2026steerable} show that early hidden states predict the outcome
of steering; we test the causal version. We steer only the prompt positions, only the
first $k$ generated tokens, or only tokens after an unsteered prefix of length $k$,
and measure the switching probability as a function of $k$. This separates three
accounts: a reply's mode is fixed by how the prompt is represented, by its first few
tokens, or accumulates throughout generation.

\subsection{Q4: Regimes and Scale}
\label{sec:q4}
We measure $\mu$, $S$, and the strength at which coherence breaks across Qwen3 models
from 0.6B to 14B, and test whether the usable window between switching and breakdown
widens with size. We also test a two-regime account of the side-effect conflict: below
the coherence break, side effects are specific to the steering vector, while above it
any push drifts the model toward a common default profile.

\fillin{Operational definition of the ``coherence break'' (what coherence score
counts as broken), and how the two-regime account is tested.}
\upd{The coherence break of a trait is the strength at which the mean coherence
score of sampled replies falls below 5, interpolated linearly in log strength
between the two neighbouring grid values. The usable window is the log distance
from the median threshold to the break, $\log a_{\mathrm{break}} - \mu$. For the
two-regime test we steer each trait vector and three norm-matched random
directions on 40 neutral prompts, at one strength below and one above the break,
and score the replies with all ten rubrics. A direction's side-effect profile is
the change from the unsteered scores on the other traits. The account predicts
that profiles of different directions, random ones included, correlate weakly
below the break and strongly above it.}

%------------------------------------------------------------
\section{Experiments}

Each question above reduces to a quantity with a confidence interval, so evaluation
for us will mean checking that these quantities are stable and not artifacts of the
judge or decoding.

\subsection{Data}
Trait prompts come from the persona-vectors release \cite{chen2025persona}
\upd{for the five traits it covers (hallucinating, sycophantic, humorous,
optimistic, impolite; used as extraction set A)}, so our
trait definitions will match a published baseline. Additionally, we will use
independently-generated held-out sets of 120 evaluation prompts per trait, so that
vectors are never evaluated on the prompts that produced them.

\fillin{How the held-out prompts were generated (which model, what instructions, any
filtering).}
\upd{All remaining extraction sets and all held-out sets were written by Claude
(Anthropic) from a fixed written brief, which is in our repository. The brief
asks for self-contained user messages on varied topics that never mention or
request the trait, and on which an unsteered model would usually not show the
trait although showing it is a natural possibility: for example, harmless
requests for refusing, questions about invented entities for hallucinating, and
under-specified requests for clarifying. An automatic check enforces the set
sizes and rejects any question that appears twice within a trait, across
extraction and held-out sets. Prompts have not been filtered by a model's
unsteered behavior; the baseline rate $g$ absorbs prompts that express the trait
without steering.} \pending{unsteered trait rate per trait, to confirm the
held-out prompts start below threshold.}

\subsection{Baselines and Controls}
\fillin{This section does not exist in the proposal and the assignment requires it.
State each baseline and what the comparison tells you.}
\upd{We use four baselines and controls.
\textbf{Unsteered model.} Every grid includes $a=0$, which gives the base rate of
each trait and the baseline coherence, and identifies $g$.
\textbf{Population-level fit.} A single sigmoid fitted to the pooled response
curve is the description used in prior work. It recovers only $w^2 + \sigma_t^2$,
so it is the baseline the hierarchical model is compared against.
\textbf{Greedy decoding.} Each prompt also gets one greedy reply per strength, so
it has a single flip point. The threshold lies between two neighbouring grid
strengths, and fitting a normal distribution to these intervals estimates
$\sigma_t$ without the hierarchical model.
\textbf{Random directions.} Three norm-matched random directions are steered at
the same strengths as the trait vectors in the side-effect runs. They show what
any push of that size does, as opposed to a push along the trait vector.}

\subsection{Evaluation Metrics}
\textbf{Reliability.} Every quantity will be estimated from at least three independent
vector sets, with prompt-bootstrap confidence intervals. Thus, we interpret
side-effect estimates only if they repeat across sets (test-retest $r > 0.7$).

\textbf{Switching.} Q1 and Q2 are scored by $S$ with its bootstrap interval per trait.
Because about ten traits give little power for a group comparison, we plan to test the
decision versus style split with a permutation test over trait labels. Greedy decoding
will give us an independent check, where each prompt then flips deterministically at
one strength, and the spread of these flip points estimates $\sigma_t$ without the
hierarchical model.

\textbf{Commitment and scale.} Q3 reports switching probability against $k$ for each
intervention. Q4 fits $\mu$ and the coherence window against log model size with
bootstrap intervals.

\textbf{Judge.} The primary judge will be Mistral-Small-24B, outside the steered model
family, scored by its expected rating, which is the probability-weighted mean over its
next-token distribution on the digits 0 to 9. Qwen3-14B will re-score a fixed 10\%
subset, and traits whose two judges disagree (meaning they have a Spearman correlation
below 0.8) will be reported but excluded from the Q2 test.

\fillin{The judge prompt (or a pointer to an appendix), and any human validation of
the judge.}
\upd{The judge prompts are given in Appendix~\ref{app:judge}. The judge reads the
rubric and one reply, and we take its next-token distribution in a single forward
pass; a score is discarded if less than 25\% of the probability falls on digits.}
\pending{human validation: agreement between judge scores and our own reading of
50 steered replies for each of three traits, and the Spearman correlation between
the two judges per trait.}

\subsection{Implementation Details}
We propose using the Qwen3 family \cite{yang2025qwen3} at 0.6B, 1.7B, 4B, 8B, and 14B
parameters because its dense models share a tokenizer and training recipe. We will use
a second family, Gemma 3 (1B, 4B, 12B, 27B), to validate the main results. The main
grid is about 96k replies per model and vector set (10 traits, 120 prompts, 10
strengths, 8 samples). We estimate 250 GPU-hours for the full study, using H100s on
the MIT Engaging cluster or RunPod.
\pending{measured throughput from the pilot, to confirm the GPU-hour estimate. If
it is exceeded, models other than Qwen3-4B use 60 prompts per trait, and Gemma 3
is limited to 4B and 12B.}

\fillin{Citation for Gemma 3.}
\upd{[Gemma Team, ``Gemma 3 Technical Report'', 2025, arXiv:2503.19786 -- check the
number and add to references.bib.]}
\fillin{Software used for steering and judging. The updated proposal dropped this.}
\upd{Steering uses forward hooks in PyTorch and Hugging Face Transformers. Judging
uses the same libraries, not vLLM, since it needs one forward pass per reply and
no generation. Every model, including the 24B judge, fits on one 80GB GPU. All
stages write append-only files and resume after interruption.}
\fillin{The 10 strength values, sampling temperature, max generation length, and
random seeds.}
\upd{Replies are sampled at temperature 1 with no nucleus or top-$k$ truncation and
are capped at 200 new tokens; Qwen3's thinking mode is disabled. Each prompt and
strength gets 8 sampled replies and one greedy reply. The grid has ten log-spaced
positive strengths plus $a=0$. Generation uses seed 0; the bootstrap and the
random directions use fixed seeds.} \pending{the ten strength values, which are
chosen from a pilot sweep on Qwen3-4B so that the grid runs from no effect to
past the coherence break.}

%------------------------------------------------------------
\section{Preliminary Results}

We have implemented the extraction, steering, and judging pipeline, and early runs
show that independently extracted vectors for the same trait agree closely.
\fillin{The number behind ``agree closely'' (cosine similarity or correlation).}
\pending{cosine similarity between vector sets A, B, and C per trait, from the
re-implemented pipeline.}

\upd{\textbf{Estimator check.} Before spending GPU time we simulated replies from
the model at the real design size (120 prompts, 10 strengths, 8 samples, baseline
rate 0.05) and refitted them, 20 times per setting. The estimator recovers $S$ in
all three settings: a switch with true $S = 0.966$ gives $0.964 \pm 0.007$ (mean
and standard deviation of the estimate), a dial with true $S = 0.034$ gives
$0.031 \pm 0.008$, and a mixed case with true $S = 0.5$ gives $0.485 \pm 0.041$.
The switch and the dial have the same $w^2 + \sigma_t^2$, so their pooled response
curves are the same and the population-level baseline cannot tell them apart.
The design therefore has enough power to separate the two pictures.

\textbf{Pipeline check.} The full pipeline runs end to end on a laptop with
Qwen3-0.6B as the steered model and Qwen3-1.7B standing in for the judge, for
one trait (humorous) and 5 prompts. This is a test of the code, not a result:
the models and sample are too small to interpret. It does show the expected
shape. The mean trait score rises from 1.2 unsteered to 5.8 at $a = 0.6$ and
falls to 2.6 at $a = 2$, where replies have degraded into repeated tokens, which
is why the fit is restricted to strengths below the coherence break.}

\fillin{Whatever results you have so far, even for one trait on Qwen3-4B, with an
interpretation relative to the baselines above.}
\pending{response curves (per-prompt and pooled) for one decision trait and one
style trait, and fitted $\mu$, $w$, $\sigma_t$, $S$ per trait on Qwen3-4B,
compared with the population-level fit and the greedy estimate of $\sigma_t$.}

\begin{table}[t]
\centering
\begin{tabular}{lcccc}
\toprule
Trait & $\mu$ & $w$ & $\sigma_t$ & $S$ \\
\midrule
\fillin{trait 1} & \pending{--} & \pending{--} & \pending{--} & \pending{--} \\
\fillin{trait 2} & \pending{--} & \pending{--} & \pending{--} & \pending{--} \\
\bottomrule
\end{tabular}
\caption{\fillin{caption} Fitted parameters per trait on Qwen3-4B.}
\label{tab:results}
\end{table}

%------------------------------------------------------------
\section{Impact Statement}
\fillin{Nothing in the proposal covers this; it is on the rubric.}
\upd{This work is meant to make steering-based safeguards more predictable. If
traits respond as switches, a strength tuned on test prompts will fully change
some new prompts and leave others untouched, so safety tools built on steering
can fail without warning; measuring $S$ shows when that risk applies. The same
vectors can be misused: a vector that suppresses a trait also amplifies it, and
our trait set includes fabrication and deception. We steer only open-weight
models that anyone can already modify, use harmless prompts, and study the shape
of the response and not ways to make harmful steering more effective. Every
measurement depends on an LLM judge, whose biases carry into our results; we
address this in part with a second judge from a different model family and a
manual check, and we exclude traits on which the judges disagree. All prompts
are synthetic, written by a language model or taken from a public release, and
no human-subject data is used.}

%------------------------------------------------------------
\section{Conclusion and Next Steps}
\fillin{2--3 sentence summary of what is done so far.}
\upd{The pipeline for all four questions is implemented and tested: extraction,
position-restricted steering, local judging, the hierarchical fit, and the
analyses for Q2 to Q4. Trait definitions, groups, and prompt sets for ten traits
are fixed. In simulation the estimator recovers $S$ at our design size. No
measurement on the target models has been run yet with this pipeline.}

Timeline: weeks 1 to 2, switching and $S$ on Qwen3-4B; week 3, commitment experiments;
weeks 4 to 5, the size ladder and second family; week 6, regime analysis and judge
cross-checks; week 7, writing. Each of us will own one question (Q1 to Q4). If traits
respond as dials, the size and regime results still stand.
\upd{Status on 19 November: the pipeline and prompt sets are complete. Remaining:
the pilot and the main grid on Qwen3-4B with the Q1 fit and Q2 test (by 2
December); commitment runs and the launch of the size ladder (by 9 December); Q4
analysis, the side-effect matrix, judge cross-checks, and Gemma 3 if budget
allows (by 16 December); then figures and writing.}

%------------------------------------------------------------
\section*{Reflection}
\fillin{Minimum 150 words, at least one Description and one Evaluation question,
with specific moments from your team's work.}
\rec{Material you could draw on: Description -- how you decided what to cut between
proposal versions (the reference list went from 16 to 9 sources); the split where
each person owns one question; tools used (Overleaf, the Engaging cluster).
Evaluation -- what worked or did not in that split, and what you would change.}

\bibliographystyle{IEEEtran}
\bibliography{references}

\appendix
\section{Judge Prompts}
\label{app:judge}
\upd{Each trait rubric inserts the trait name and its one-paragraph description
(listed in \texttt{configs/traits.yaml} in our repository), the user prompt, and
the reply:

\begin{quote}\small
You are evaluating whether the model's response displays the trait:
**\{trait\}**. \{description\}

Prompt: [QUESTION START] \{question\} [QUESTION END]

Model response: [ANSWER START] \{answer\} [ANSWER END]

Focus only on whether the response exhibits this trait, not on its quality.

Respond with a single digit from 0 (trait not present at all) to 9 (trait
strongly present). Don't say anything else, just the digit.
\end{quote}

The coherence rubric asks for a single digit, where 0 means the answer makes no
sense at all (gibberish, endless repetition, or not valid language), 5 means it
is understandable but rambling, confused, or drifting away from the prompt, and
9 means clear, fluent language that addresses the prompt. It states that an
answer can be rude, wrong, strange in style, or a refusal and still be fully
coherent.}

\end{document}